\chapter{Numerical analysis and reproducible scientific computation}\label{ch:numerical}
\begin{readergoal}Translate mathematical equations into algorithms without overlooking finite precision, memory layout, and validation.\end{readergoal}
\section{A formula is not yet an algorithm}
A mathematical expression can define an exact quantity while computer code approximates it. Floating-point arithmetic rounds values and is not always associative: $(a+b)+c$ and $a+(b+c)$ can differ. An apparently exact bin count computed by $\lfloor T/\Delta\rfloor$ can be off by one if $T/\Delta$ is represented just below an integer. Implementations must state endpoint conventions and test boundary cases rather than rely on a universal ad hoc epsilon.

\section{Conditioning and algorithmic stability}
A well-posed problem has a solution that exists, is unique, and depends continuously on the data in an appropriate sense. An ill-conditioned problem may magnify small perturbations even under perfect implementation. A numerically unstable algorithm may magnify errors beyond what conditioning requires. Computing least-squares coefficients via QR or SVD is typically more stable than directly inverting $X^\top X$. Regularization stabilizes some inverse problems by choosing a constrained or penalized solution, at the cost of explicit modeling assumptions.

\section{Memory and complexity}
A dense sample covariance among $N$ neurons requires $O(N^2)$ storage, while a full matrix decomposition can be computationally expensive for large $N$. When sample count $T$ is much smaller than $N$, methods operating on a $T\times T$ Gram matrix can be more efficient. Signal operations can frequently be evaluated chunk by chunk, although filtering chunks requires state continuity or overlap handling. Memory mapping avoids loading full arrays but does not eliminate I/O costs.

\section{Testing the mathematical contract}
Good numerical tests use known synthetic inputs: a sinusoid with known frequency for PSD estimation; a complex exponential for phase operations; a Poisson simulation for count statistics; a low-rank signal with controlled noise for PCA; and a known VAR for Granger analysis. Property-based tests can verify constraints such as covariance symmetry, nonnegative variance, $\mathrm{PLV}\in[0,1]$, and consistency of concatenated block processing with full-array processing within a stated tolerance.

\section{Scientific metadata as part of the result}
An estimate needs sampling rate, time origin, measurement units, reference configuration, filtering parameters, electrode and trial identifiers, quality masks, random seeds, and software versions. Metadata are part of the mathematical map between the array and its real-world interpretation. A floating-point matrix detached from its coordinates can be mathematically valid yet scientifically unusable.

\begin{exercises}
\item Explain why zero-padding an FFT does not improve physical spectral resolution.
\item Design a test for chunked FIR filtering that detects boundary discontinuities.
\item List the metadata necessary to interpret a Hilbert-envelope moment series from 10-second windows.
\end{exercises}